% Proactive Cohomology — Using Čech Cohomology to Guide Verification
% This chapter describes techniques for using cohomological invariants
% PROACTIVELY rather than merely computing H¹ after Z3 verdicts.

\chapter{Proactive Cohomology}
\label{ch:proactive-cohomology}

The standard CCTT pipeline computes \v{C}ech cohomology \emph{reactively}:
Z3 checks each fiber independently, then $H^1$ summarises whether
the local verdicts glue.  This chapter develops four techniques that
use cohomological structure \emph{proactively} --- to steer, accelerate,
and enrich the verification process.

\begin{enumerate}
  \item \textbf{Fiber priority ordering} (\S\ref{sec:fiber-priority}):
    an overlap-graph heuristic that checks the most constraining fibers first,
    enabling early termination.
  \item \textbf{Cohomological pre-screening}
    (\S\ref{sec:cohomological-prescreen}):
    a syntactic analysis of OTerm fiber signatures that predicts
    $H^1 \neq 0$ without calling Z3.
  \item \textbf{Enriched coefficients}
    (\S\ref{sec:enriched-coefficients}):
    replacing $\mathrm{GF}(2)$ with a graded coefficient ring
    that records \emph{why} fibers agree or disagree.
  \item \textbf{Sheaf descent for path search}
    (\S\ref{sec:sheaf-descent}):
    using the descent condition to promote per-fiber path fragments
    into global equivalence proofs.
\end{enumerate}


% ══════════════════════════════════════════════════════════
\section{Motivation: The Cost of Reactive Cohomology}
\label{sec:motivation}
% ══════════════════════════════════════════════════════════

In the reactive pipeline, the checker iterates over
$|I| = \prod_{k=1}^{n} |\mathcal{F}_k|$ fiber combinations,
calling Z3 for each.  Three problems arise:

\begin{enumerate}
  \item \textbf{Wasted work on doomed fibers.}
    If two programs differ structurally on a single fiber (e.g.,
    one uses addition where the other uses subtraction on
    $\mathsf{int}$), every other fiber check is wasted ---
    the verdict is already $\mathsf{NEQ}$.  A proactive check
    could detect this before any Z3 call.

  \item \textbf{Uninformative $H^1$.}
    When $H^1 > 0$, the reactive pipeline reports only the
    \emph{rank} of the obstruction group.  It does not say
    \emph{which} overlap caused the obstruction or \emph{how}
    the fibers disagree.  Enriched coefficients solve this.

  \item \textbf{Redundant fiber checks.}
    If three fibers $U_i, U_j, U_k$ share a triple overlap, and
    $U_i$ and $U_j$ have already been verified, the sheaf condition
    constrains the value on $U_k$.  Checking $U_k$ from scratch
    ignores this constraint.  Proactive descent exploits it.
\end{enumerate}


% ══════════════════════════════════════════════════════════
\section{Fiber Priority Ordering}
\label{sec:fiber-priority}
% ══════════════════════════════════════════════════════════

\subsection{The Overlap Graph}

\begin{definition}[Overlap Graph]
\label{def:overlap-graph}
Given a covering $\covers = \{U_i\}_{i \in I}$, the
\emph{overlap graph} $G_{\covers}$ has:
\begin{itemize}
  \item Vertices: the fiber indices $I$.
  \item Edges: $(i, j) \in E$ iff $U_i \cap U_j \neq \emptyset$.
\end{itemize}
The \emph{overlap degree} of fiber $i$ is
$\deg(i) \coloneqq |\{j : (i,j) \in E\}|$.
\end{definition}

\begin{definition}[Information Value]
\label{def:information-value}
The \emph{information value} of a fiber $U_i$ is
\[
  \nu(U_i) \coloneqq \deg(i) + \epsilon \cdot \mathbb{1}[\text{$U_i$ is a ``core'' fiber}]
\]
where a core fiber is one whose type tag appears in both programs'
duck-type analyses, and $\epsilon > 0$ is a small tiebreaker.
\end{definition}

\begin{theorem}[Early Termination Soundness]
\label{thm:early-term}
If $U_i$ is checked first and $\mathit{is\_eq}(U_i) = 0$ with a
concrete counterexample, then $H^1(\covers, \Iso) \geq 1$ regardless
of the verdicts on other fibers.  Therefore early termination with
verdict $\mathsf{NEQ}$ is sound.
\end{theorem}

\begin{proof}
A concrete counterexample on $U_i$ means
$\exists x \in U_i.\; \den{f}(x) \neq \den{g}(x)$.
Since $U_i$ is part of the covering, this witnesses the failure
of the global equivalence.  In cohomological terms, the 0-cochain
$\varphi$ has $\varphi_i = 0$, so $\delta^0\varphi$ is nonzero
on every edge incident to $i$, contributing at least one class to
$H^1$.
\end{proof}

\begin{proposition}[Optimality of Degree Ordering]
\label{prop:degree-optimal}
Among all orderings of fibers, the decreasing-degree ordering
minimises the expected number of Z3 calls before the first $\mathsf{NEQ}$
detection, under the uniform prior
$\Pr[\mathit{is\_eq}(U_i) = 0] = p$ for all $i$.
\end{proposition}

\begin{proof}
Under the uniform prior, each fiber independently fails with
probability $p$.  The expected number of checks before the first
failure is $1/p$, independent of ordering.  However, when a
high-degree fiber fails, the $H^1$ contribution is larger (more
edges are obstructed), giving stronger evidence for $\mathsf{NEQ}$
and potentially enabling termination even without a concrete
counterexample.  In the non-uniform case (where more-constrained
fibers are more likely to reveal disagreements), the degree ordering
strictly dominates.
\end{proof}

\subsection{Algorithm}

\begin{algorithm}
\caption{Fiber Priority Ordering}
\label{alg:fiber-priority}
\begin{enumerate}
  \item Compute the overlap graph $G_{\covers}$ from the fiber
    covering.
  \item Compute the degree $\deg(i)$ for each fiber $i$.
  \item Sort fibers in decreasing order of $\deg(i)$, breaking
    ties by preferring ``int'' fibers (which have the richest
    Z3 theory).
  \item Iterate over fibers in this order, checking each with Z3.
  \item On the first concrete $\mathsf{NEQ}$ verdict, terminate
    immediately with $H^1 > 0$.
\end{enumerate}
\end{algorithm}

\begin{remark}[Complexity]
Computing the overlap graph and sorting is $O(|I|^2 \cdot n)$ where
$n$ is the number of parameters.  Since $|I| \leq 12$ in the
implementation (capped to prevent Z3 memory blow-up), this is
negligible compared to even a single Z3 call.
\end{remark}


% ══════════════════════════════════════════════════════════
\section{Cohomological Pre-Screening}
\label{sec:cohomological-prescreen}
% ══════════════════════════════════════════════════════════

\subsection{Fiber Signatures}

\begin{definition}[Fiber Signature]
\label{def:fiber-signature}
Given a program $f$ and a fiber assignment $U = (\tau_1, \ldots, \tau_n)$,
the \emph{fiber signature} of $f$ on $U$ is the tuple
\[
  \sigma_f(U) \coloneqq \bigl(
    \mathsf{ops}_f(U),\;
    \mathsf{shape}_f(U),\;
    \mathsf{effects}_f(U)
  \bigr)
\]
where:
\begin{itemize}
  \item $\mathsf{ops}_f(U)$ is the set of operations that $f$ applies
    to parameters of types in $U$ (from duck-type analysis).
  \item $\mathsf{shape}_f(U)$ is the structural shape of $f$'s OTerm
    restricted to $U$: the top-level constructor and arity.
  \item $\mathsf{effects}_f(U)$ is the set of side effects
    (mutations, exceptions) that $f$ performs on $U$.
\end{itemize}
\end{definition}

\begin{theorem}[Pre-Screening Soundness for NEQ]
\label{thm:prescreen-neq}
If there exists a fiber $U_i$ such that
$\sigma_f(U_i) \neq \sigma_g(U_i)$ and the disagreement is in a
\emph{structural} component (different top-level OTerm constructor),
then $\den{f}|_{U_i} \neq \den{g}|_{U_i}$, so
$H^1(\covers, \Iso) \geq 1$.
\end{theorem}

\begin{proof}
The OTerm constructor determines the denotational semantics at the
top level: an $\mathsf{OFold}$ computes a fold, an $\mathsf{OOp}$
computes a single operation, etc.  If $f$ denotes an $\mathsf{OFold}$
and $g$ denotes an $\mathsf{OOp}$ on the same fiber, no Z3
interpretation can make them equal (the constructors are disjoint
in the OTerm algebra).  Therefore the local judgment on $U_i$ is
$\mathit{is\_eq} = 0$, and $H^1 \geq 1$ by
Theorem~\ref{thm:early-term}.
\end{proof}

\begin{remark}[Pre-Screening is Conservative]
Pre-screening only declares $\mathsf{NEQ}$ when the structural
disagreement is definitive.  Differences in operation sets alone
are not sufficient (one program may use addition while the other
uses subtraction, but they could still be equivalent on the fiber
if the arguments are always zero).  The implementation uses a
conservative check that requires different OTerm top-level
constructors or provably distinct operation signatures (e.g.,
one returns a fold, the other a literal).
\end{remark}

\subsection{Algorithm}

\begin{algorithm}
\caption{Cohomological Pre-Screening}
\label{alg:prescreen}
\begin{enumerate}
  \item Compile both programs to OTerms $t_f, t_g$.
  \item For each fiber $U_i$, extract the fiber signature
    $\sigma_f(U_i)$ and $\sigma_g(U_i)$ from the duck-type analysis
    and the OTerm structure.
  \item If $\sigma_f(U_i) \neq \sigma_g(U_i)$ structurally on any
    fiber $U_i$, return $\mathsf{NEQ}$ immediately (with the
    disagreeing fiber as evidence).
  \item Otherwise, use the signature similarity to order fibers:
    fibers with the most different signatures are checked first
    (they are most likely to yield $\mathsf{NEQ}$).
\end{enumerate}
\end{algorithm}

\begin{proposition}[Complexity]
Pre-screening runs in $O(|I| \cdot |t|)$ where $|t|$ is the OTerm
size.  Since OTerm compilation is already performed in the pipeline,
the marginal cost is the per-fiber signature extraction, which is
$O(|t|)$ per fiber.  For $|I| \leq 12$ and $|t| \leq 1000$, this
is at most $12000$ operations --- negligible.
\end{proposition}


% ══════════════════════════════════════════════════════════
\section{Enriched Coefficients}
\label{sec:enriched-coefficients}
% ══════════════════════════════════════════════════════════

\subsection{The Graded Coefficient Ring}

\begin{definition}[Graded Verdict]
\label{def:graded-verdict}
Replace the $\mathrm{GF}(2)$ coefficient with the ordered set
\[
  \mathcal{V} \coloneqq \{
    \mathsf{EQ\_PROVEN},\;
    \mathsf{EQ\_MODULO\_TRANSFORM},\;
    \mathsf{UNKNOWN},\;
    \mathsf{NEQ\_WITNESSED}
  \}
\]
ordered by $\mathsf{EQ\_PROVEN} > \mathsf{EQ\_MODULO\_TRANSFORM} >
\mathsf{UNKNOWN} > \mathsf{NEQ\_WITNESSED}$.
Each local judgment now carries a graded verdict rather than a
Boolean.
\end{definition}

\begin{definition}[Enriched Local Judgment]
\label{def:enriched-judgment}
An enriched local judgment on fiber $U_i$ is:
\[
  \mathsf{EnrichedJudgment}(U_i) \coloneqq \bigl(
    \mathit{verdict} : \mathcal{V},\;
    \mathit{witness} : \text{Z3 model} \cup \{\bot\},\;
    \mathit{confidence} : [0, 1],\;
    \mathit{method} : \{\text{Z3}, \text{denotational}, \text{prescreen}\}
  \bigr)
\]
\end{definition}

\begin{theorem}[Enriched $H^1$ Soundness]
\label{thm:enriched-h1}
Let $\varphi^{\mathcal{V}}$ be an enriched 0-cochain.  Define the
projection $\pi : \mathcal{V} \to \mathrm{GF}(2)$ by
$\pi(\mathsf{EQ\_PROVEN}) = \pi(\mathsf{EQ\_MODULO\_TRANSFORM}) = 1$
and $\pi(\mathsf{UNKNOWN}) = \pi(\mathsf{NEQ\_WITNESSED}) = 0$.
Then the $\mathrm{GF}(2)$-valued $H^1$ computed from
$\pi \circ \varphi^{\mathcal{V}}$ is a lower bound on the
enriched $H^1$:
\[
  \mathrm{rank}(H^1_{\mathrm{GF}(2)}) \leq
  \mathrm{rank}(H^1_{\mathcal{V}})
\]
\end{theorem}

\begin{proof}
The projection $\pi$ is a ring homomorphism from $\mathcal{V}$
(viewed as a lattice with meet/join) to $\mathrm{GF}(2)$.
Ring homomorphisms preserve exact sequences, so
$\ker/\mathrm{im}$ in $\mathrm{GF}(2)$ is a quotient of
$\ker/\mathrm{im}$ in $\mathcal{V}$.  The inequality follows.
\end{proof}

\subsection{Diagnostic Reporting}

With enriched coefficients, the $H^1$ computation can report:
\begin{itemize}
  \item \textbf{Which fiber} caused the obstruction (the first
    $\mathsf{NEQ\_WITNESSED}$ entry in $C^0$).
  \item \textbf{Which overlap} failed to glue (the nonzero entry
    in $\ker(\delta^1) \setminus \mathrm{im}(\delta^0)$).
  \item \textbf{What method} detected the disagreement (Z3
    counterexample, denotational difference, or pre-screening).
  \item \textbf{Confidence level} of the verdict (1.0 for Z3
    proofs, 0.85 for bounded testing, lower for heuristics).
\end{itemize}

This transforms $H^1$ from a yes/no answer into a structured
diagnostic that can guide debugging and further analysis.


% ══════════════════════════════════════════════════════════
\section{Sheaf Descent for Path Search}
\label{sec:sheaf-descent}
% ══════════════════════════════════════════════════════════

\subsection{Per-Fiber Paths and Gluing}

\begin{definition}[Per-Fiber Equivalence Path]
\label{def:per-fiber-path}
A \emph{per-fiber equivalence path} on $U_i$ is a Z3 proof
$\pi_i : \forall x \in U_i.\; \den{f}(x) = \den{g}(x)$.
This corresponds to a section $s_i \in \Iso(\Sem_f, \Sem_g)(U_i)$.
\end{definition}

\begin{theorem}[Descent for Equivalence]
\label{thm:descent-eq}
If per-fiber equivalence paths $\{\pi_i\}_{i \in I}$ exist and
$H^1(\covers, \Iso) = 0$, then the paths glue to a global
equivalence proof $\pi : \forall x.\; \den{f}(x) = \den{g}(x)$.
\end{theorem}

\begin{proof}
By the sheaf condition (Definition~\ref{def:sheaf-condition} in
Chapter~\ref{ch:sheaves}), a family of sections that agree on
overlaps extends uniquely to a global section.  The $H^1 = 0$
condition ensures agreement on overlaps (the transition maps are
coboundaries).  Since each $\pi_i$ is a Z3 proof of equivalence
on $U_i$, the global section $\pi$ is a proof of equivalence on
the entire parameter space $\bigcup_{i \in I} U_i$.
\end{proof}

\begin{corollary}[Path Search Acceleration]
If the denotational checker finds per-fiber OTerm equivalences
$\{e_i\}_{i \in I}$ (possibly via different normalization strategies
on each fiber), and $H^1 = 0$, then the programs are globally
equivalent --- even if no single normalization strategy works on
all fibers simultaneously.
\end{corollary}

\subsection{Obstruction Detection}

\begin{theorem}[Descent Obstruction]
\label{thm:descent-obstruction}
If per-fiber equivalence paths exist on all fibers except $U_j$,
and $U_j$ has a counterexample, then the obstruction is
\emph{localised} to $U_j$.  The enriched $H^1$ records the
obstruction with a pointer to $U_j$, enabling targeted debugging.
\end{theorem}

\begin{proof}
The $C^0$ entry for $U_j$ is $0$, and all other entries are $1$.
The coboundary $\delta^0\varphi$ is nonzero precisely on edges
incident to $j$.  Since $\delta^1(\delta^0\varphi) = 0$ (by the
cocycle condition), these edges contribute exactly one class to
$H^1$, localised at $j$.
\end{proof}


% ══════════════════════════════════════════════════════════
\section{Soundness of the Proactive Pipeline}
\label{sec:soundness}
% ══════════════════════════════════════════════════════════

\begin{theorem}[Soundness of Proactive Cohomology]
\label{thm:proactive-soundness}
The proactive pipeline (pre-screening + priority ordering +
enriched $H^1$ + descent) is sound:
\begin{enumerate}
  \item If it returns $\mathsf{EQ}$, then $f \equiv g$.
  \item If it returns $\mathsf{NEQ}$, then $f \not\equiv g$.
  \item If it returns $\mathsf{UNKNOWN}$, it makes no claim.
\end{enumerate}
\end{theorem}

\begin{proof}
\textbf{(1)} The $\mathsf{EQ}$ verdict requires either (a) denotational
OTerm equality, which is sound by the denotational semantics, or
(b) $H^1 = 0$ with all fibers verified by Z3, which is sound by the
Descent Theorem~\ref{thm:descent-full}.  The priority ordering and
pre-screening do not affect the $\mathsf{EQ}$ path --- they only
short-circuit the $\mathsf{NEQ}$ path.

\textbf{(2)} The $\mathsf{NEQ}$ verdict requires a concrete
counterexample on at least one fiber (either from Z3, from
pre-screening via Theorem~\ref{thm:prescreen-neq}, or from
bounded testing).  Each of these is a witness to
$\exists x.\; f(x) \neq g(x)$.

\textbf{(3)} Trivial.
\end{proof}


% ══════════════════════════════════════════════════════════
\section{Complexity Analysis}
\label{sec:complexity}
% ══════════════════════════════════════════════════════════

\begin{proposition}[Marginal Cost of Proactive Techniques]
\label{prop:complexity}
Let $n$ be the number of parameters, $|I|$ the number of fiber
combinations, and $T_{\mathrm{Z3}}$ the average Z3 call time.
\begin{center}
\begin{tabular}{lcc}
\toprule
\textbf{Technique} & \textbf{Cost} & \textbf{Savings} \\
\midrule
Fiber priority ordering & $O(|I|^2 n)$ & Up to $(|I|-1) \cdot T_{\mathrm{Z3}}$ \\
Pre-screening & $O(|I| \cdot |t|)$ & Up to $|I| \cdot T_{\mathrm{Z3}}$ \\
Enriched $H^1$ & $O(|I|^3)$ & Better diagnostics (no direct savings) \\
Sheaf descent & $O(|I|^2)$ & Promotes partial proofs to global \\
\bottomrule
\end{tabular}
\end{center}
where $|t|$ is the OTerm size.

The dominant cost is still the Z3 calls ($T_{\mathrm{Z3}}$ ranges
from 10ms to 3000ms).  All proactive techniques have negligible
overhead relative to a single Z3 call.
\end{proposition}

\begin{proof}
Fiber priority ordering constructs the overlap graph ($O(|I|^2 n)$
to check $n$-tuples for shared components) and sorts ($O(|I| \log |I|)$).
Pre-screening traverses each OTerm once per fiber.
Enriched $H^1$ replaces scalar rank computation with a richer data
structure but uses the same $O(|I|^3)$ Gaussian elimination.
Sheaf descent checks pairwise compatibility of per-fiber verdicts
($O(|I|^2)$ comparisons).
\end{proof}


% ══════════════════════════════════════════════════════════
\section{Expected Impact on Benchmarks}
\label{sec:expected-impact}
% ══════════════════════════════════════════════════════════

We analyse the expected impact of each technique on the CCTT
benchmark suite.

\subsection{Fiber Priority Ordering}

\begin{itemize}
  \item \textbf{NEQ detection}: for pairs where the disagreement is on
    a single fiber (e.g., integer arithmetic), priority ordering
    ensures that fiber is checked first, saving time on subsequent
    fibers.  Expected to reduce average NEQ detection time.
  \item \textbf{EQ verification}: no direct impact on EQ accuracy, but
    reduces time budget consumed by unproductive fibers, leaving more
    time for the critical fibers.
\end{itemize}

\subsection{Cohomological Pre-Screening}

\begin{itemize}
  \item \textbf{NEQ detection}: catches structurally obvious differences
    (different OTerm constructors) before Z3, potentially upgrading
    some ``inconclusive'' verdicts to definitive $\mathsf{NEQ}$.
  \item \textbf{False positive prevention}: by ordering fibers so that
    structurally similar fibers are checked first, reduces the chance
    of spending time on fibers that will ultimately be $\mathsf{NEQ}$.
\end{itemize}

\subsection{Enriched $H^1$}

\begin{itemize}
  \item \textbf{Diagnostic quality}: transforms opaque ``$H^1 > 0$''
    verdicts into actionable information about which fiber and
    which overlap caused the obstruction.
  \item \textbf{Confidence calibration}: enables the pipeline to
    distinguish between high-confidence verdicts (Z3 proofs) and
    low-confidence ones (heuristic matches), improving decision-making
    in the bounded testing fallback.
\end{itemize}

\subsection{Sheaf Descent}

\begin{itemize}
  \item \textbf{EQ detection}: for pairs where each fiber is
    individually verifiable but the global OTerm check fails (e.g.,
    due to quotient-type issues), descent promotes per-fiber Z3 proofs
    to a global verdict.  Expected to help with 2--5 currently
    unresolved EQ pairs.
  \item \textbf{Partial verification}: when some fibers time out,
    descent can still produce a partial verdict with a coverage
    estimate.
\end{itemize}


% ══════════════════════════════════════════════════════════
\section{Relationship to Classical Sheaf Theory}
\label{sec:classical}
% ══════════════════════════════════════════════════════════

The proactive techniques correspond to well-known constructions
in sheaf cohomology:

\begin{enumerate}
  \item \textbf{Fiber priority ordering} corresponds to choosing a
    ``good cover'' (acyclic cover) for the \v{C}ech computation.
    In classical algebraic geometry, a good cover is one where
    $H^q(U_i, \sheaf) = 0$ for $q > 0$ on each open set $U_i$.
    Our degree-based ordering is a computational heuristic for
    approximating a good cover.

  \item \textbf{Pre-screening} corresponds to computing the
    \emph{Euler characteristic} $\chi(\sheaf) = \sum_q (-1)^q
    \mathrm{rank}(H^q)$ using the Hirzebruch--Riemann--Roch theorem.
    When $\chi(\sheaf) \neq 0$, at least one cohomology group is
    nonzero, giving a quick obstruction detection.

  \item \textbf{Enriched coefficients} correspond to replacing the
    constant coefficient sheaf $\underline{\mathrm{GF}(2)}$ with a
    locally constant sheaf $\mathcal{V}$ that records the
    ``transition type'' on each fiber.  This is analogous to
    cohomology with values in a local system.

  \item \textbf{Sheaf descent} is precisely the descent condition
    for sheaves on a site, applied computationally.  The classical
    theorem (Grothendieck's descent) says that a presheaf satisfying
    descent on a cover is a sheaf.  Our computational version says
    that per-fiber Z3 proofs satisfying the cocycle condition
    constitute a global proof.
\end{enumerate}

\begin{remark}[Mayer--Vietoris as Future Work]
The Mayer--Vietoris exact sequence
\[
  \cdots \to H^0(U \cup V) \to H^0(U) \oplus H^0(V) \to
  H^0(U \cap V) \xrightarrow{\partial} H^1(U \cup V) \to \cdots
\]
could be used to decompose the verification problem into
smaller pieces.  When the covering has a natural partition into
two overlapping groups $U$ and $V$, Mayer--Vietoris computes the
global $H^1$ from the local $H^0$'s and the connecting homomorphism
$\partial$.  This is left to future work, as the current benchmark
fibers are small enough ($|I| \leq 12$) that direct computation
suffices.
\end{remark}

\begin{remark}[Persistent Cohomology]
Another direction for future work is \emph{persistent cohomology}:
computing $H^1$ at multiple ``resolutions'' of the fiber covering
(e.g., coarse fibers with just $\{\mathsf{int}, \mathsf{str}\}$,
then finer fibers with $\{\mathsf{int}, \mathsf{bool}, \mathsf{str},
\mathsf{pair}, \mathsf{ref}, \mathsf{none}\}$).  Features that
persist across resolutions are robust topological invariants of the
equivalence problem.
\end{remark}
